\chapter{Expanded Welch--Satterthwaite Method}
\label{chap:expanded}

The expanded method retains the bias-corrected difference, standard error, and
test statistic from Chapter~\ref{chap:baselines}. It changes only the component
degrees of freedom used in the Welch--Satterthwaite reference. This chapter
derives those degrees of freedom from the sampling variability of the complete
MI variance estimator.

\section{Roadmap from the Welch--Satterthwaite equation}

For independent variance estimates \(s_i^2\), positive weights \(k_i\), and
component degrees of freedom \(\nu_i\), the Welch--Satterthwaite equation is
\begin{equation}
  \nu
  =
  \frac{
  \left(\sum_i k_i s_i^2\right)^2
  }{
  \sum_i(k_i s_i^2)^2/\nu_i
  }.
  \label{eq:ws-general}
\end{equation}
For the two MI estimates, the variance contributions are
\begin{equation}
  k_Ps_P^2=\frac{\Vhat(P)}{n_P},
  \qquad
  k_Qs_Q^2=\frac{\Vhat(Q)}{n_Q}.
  \label{eq:ws-mi-components}
\end{equation}
These are the two terms in the squared standard error. The remaining quantities
are the component degrees of freedom \(\nu_P\) and \(\nu_Q\), which should
measure how reliably \(\Vhat(P)\) and \(\Vhat(Q)\) estimate their population
counterparts.

The derivation proceeds in four stages. First, Satterthwaite moment matching
expresses \(\nu_P\) in terms of the mean and variance of \(\Vhat(P)\). Second,
a one-cell perturbation gives the observation-level sensitivity of \(V(P)\).
Third, the variability of these sensitivities gives the sampling variance of
\(\Vhat(P)\). Finally, the two estimated component degrees of freedom are
combined through Equation~\eqref{eq:ws-general}.

\section{Component degrees of freedom by moment matching}

Satterthwaite represents a positive variance estimate by a scaled chi-squared
variable. For the first population, write the working model as
\begin{equation}
  \Vhat(P)
  \ \dot\sim\
  \E\{\Vhat(P)\}
  \frac{\chi^2_{\nu_P}}{\nu_P}.
  \label{eq:scaled-chi-working}
\end{equation}
The model is a moment-matching approximation. It does not assert that the
finite-sample distribution of \(\Vhat(P)\) is exactly chi-squared.

For \(X\sim\chi^2_{\nu_P}\),
\begin{equation}
  \E(X)=\nu_P,
  \qquad
  \Var(X)=2\nu_P.
\end{equation}
The scaling in Equation~\eqref{eq:scaled-chi-working} therefore matches the
mean automatically:
\begin{equation}
  \E\left[
  \E\{\Vhat(P)\}\frac{X}{\nu_P}
  \right]
  =\E\{\Vhat(P)\}.
\end{equation}
Its variance is
\begin{align}
  \Var\left[
  \E\{\Vhat(P)\}\frac{X}{\nu_P}
  \right]
  &=
  \frac{[\E\{\Vhat(P)\}]^2}{\nu_P^2}
  \Var(X)\\
  &=
  \frac{2[\E\{\Vhat(P)\}]^2}{\nu_P}.
\end{align}
Matching this quantity to \(\Var\{\Vhat(P)\}\) gives
\begin{equation}
  \boxed{
  \nu_P
  =
  \frac{2[\E\{\Vhat(P)\}]^2}
       {\Var\{\Vhat(P)\}}
  }.
  \label{eq:component-df-moments}
\end{equation}
To first order, \(\E\{\Vhat(P)\}\approx V(P)\). The required calculation is
therefore the sampling variance of \(\Vhat(P)\).

\section{A one-cell perturbation of the distribution}

The sampling behaviour of a smooth plug-in functional can be obtained by
measuring its first-order response to one observation. Fix a cell
\((x,y)\) and define the perturbed distribution
\begin{equation}
  P_\varepsilon
  =(1-\varepsilon)P+\varepsilon\delta_{(x,y)},
  \label{eq:contamination-path}
\end{equation}
where \(\delta_{(x,y)}\) places probability one on cell \((x,y)\). At
\(\varepsilon=0\), the distribution is \(P\). Increasing \(\varepsilon\)
moves a small amount of probability toward the selected cell while preserving
total probability one.

For an arbitrary cell \((i,j)\),
\begin{equation}
  p_{ij}(\varepsilon)
  =(1-\varepsilon)p_{ij}
  +\varepsilon\mathbf 1\{i=x,j=y\}.
\end{equation}
The joint and marginal probability derivatives at the original distribution
are
\begin{align}
  \left.\frac{\mathrm d p_{ij}(\varepsilon)}{\mathrm d\varepsilon}
  \right|_{0}
  &=\mathbf 1\{i=x,j=y\}-p_{ij},
  \label{eq:joint-perturbation-derivative}\\
  \left.\frac{\mathrm d p_{i+}(\varepsilon)}{\mathrm d\varepsilon}
  \right|_{0}
  &=\mathbf 1\{i=x\}-p_{i+},\\
  \left.\frac{\mathrm d p_{+j}(\varepsilon)}{\mathrm d\varepsilon}
  \right|_{0}
  &=\mathbf 1\{j=y\}-p_{+j}.
\end{align}
The derivatives of the joint probabilities sum to zero because
\begin{equation}
  \sum_{i,j}
  \{\mathbf 1(i=x,j=y)-p_{ij}\}
  =1-1=0.
\end{equation}

\section{Change in pointwise mutual information and MI}

Along the perturbation path, pointwise MI is
\begin{equation}
  \pmi_{P_\varepsilon}(i,j)
  =\log p_{ij}(\varepsilon)
  -\log p_{i+}(\varepsilon)
  -\log p_{+j}(\varepsilon).
\end{equation}
Differentiating the three logarithms and substituting the probability
derivatives gives
\begin{equation}
  \left.
  \frac{\mathrm d}{\mathrm d\varepsilon}
  \pmi_{P_\varepsilon}(i,j)
  \right|_{0}
  =
  \frac{\mathbf 1\{i=x,j=y\}}{p_{ij}}
  -\frac{\mathbf 1\{i=x\}}{p_{i+}}
  -\frac{\mathbf 1\{j=y\}}{p_{+j}}
  +1.
  \label{eq:pmi-derivative}
\end{equation}

MI along the path is
\begin{equation}
  I(P_\varepsilon)
  =\sum_{i,j}p_{ij}(\varepsilon)
  \pmi_{P_\varepsilon}(i,j).
\end{equation}
The product rule separates its derivative into the change in probability
weights and the change in pointwise-MI values:
\begin{align}
  \left.\frac{\mathrm d I(P_\varepsilon)}{\mathrm d\varepsilon}
  \right|_{0}
  ={}&
  \sum_{i,j}
  \{\mathbf 1(i=x,j=y)-p_{ij}\}\pmi_P(i,j)
  \notag\\
  &+
  \sum_{i,j}p_{ij}
  \left.\frac{\mathrm d\pmi_{P_\varepsilon}(i,j)}
  {\mathrm d\varepsilon}\right|_{0}.
\end{align}
The first sum is \(\pmi_P(x,y)-I(P)\). Substitution of Equation
\eqref{eq:pmi-derivative} into the second sum gives \(1-1-1+1=0\). Hence
\begin{equation}
  \boxed{
  \left.\frac{\mathrm d I(P_\varepsilon)}{\mathrm d\varepsilon}
  \right|_{0}
  =\pmi_P(x,y)-I(P)
  }.
  \label{eq:mi-derivative}
\end{equation}
This result recovers the centred pointwise-MI influence function used in
Chapter~\ref{chap:baselines}.

\section{Change in the MI variance}

The target is the first-order change in
\begin{equation}
  V(P)
  =\Var_P\{\pmi_P(X,Y)\}.
\end{equation}
Write this variance as
\begin{equation}
  V(P)=M_2(P)-I(P)^2,
  \qquad
  M_2(P)=\E_P\{\pmi_P(X,Y)^2\}.
  \label{eq:v-second-moment}
\end{equation}
Define the observation-level variance sensitivity by
\begin{equation}
  g_P(x,y)
  =
  \left.
  \frac{\mathrm d V(P_\varepsilon)}{\mathrm d\varepsilon}
  \right|_{0}.
  \label{eq:g-definition}
\end{equation}
Differentiating Equation~\eqref{eq:v-second-moment} gives
\begin{equation}
  g_P(x,y)
  =
  \left.\frac{\mathrm d M_2(P_\varepsilon)}{\mathrm d\varepsilon}
  \right|_{0}
  -2I(P)
  \left.\frac{\mathrm d I(P_\varepsilon)}{\mathrm d\varepsilon}
  \right|_{0}.
  \label{eq:g-decomposition}
\end{equation}

It remains to differentiate
\begin{equation}
  M_2(P_\varepsilon)
  =\sum_{i,j}p_{ij}(\varepsilon)
  \pmi_{P_\varepsilon}(i,j)^2.
\end{equation}
The product rule gives
\begin{align}
  \left.\frac{\mathrm d M_2(P_\varepsilon)}{\mathrm d\varepsilon}
  \right|_{0}
  ={}&
  \sum_{i,j}
  \left.\frac{\mathrm d p_{ij}(\varepsilon)}{\mathrm d\varepsilon}
  \right|_{0}
  \pmi_P(i,j)^2
  \notag\\
  &+
  2\sum_{i,j}p_{ij}\pmi_P(i,j)
  \left.\frac{\mathrm d\pmi_{P_\varepsilon}(i,j)}
  {\mathrm d\varepsilon}\right|_{0}.
  \label{eq:m2-product-rule}
\end{align}
The first sum is
\begin{equation}
  \pmi_P(x,y)^2-M_2(P).
  \label{eq:m2-weight-term}
\end{equation}

For the second sum, define the probability-weighted row and column means
\begin{align}
  \E_P\{\pmi_P(X,Y)\mid X=x\}
  &=\frac{\sum_j p_{xj}\pmi_P(x,j)}{p_{x+}},\\
  \E_P\{\pmi_P(X,Y)\mid Y=y\}
  &=\frac{\sum_i p_{iy}\pmi_P(i,y)}{p_{+y}}.
\end{align}
Substitution of Equation~\eqref{eq:pmi-derivative} gives
\begin{align}
  &\sum_{i,j}p_{ij}\pmi_P(i,j)
  \left.\frac{\mathrm d\pmi_{P_\varepsilon}(i,j)}
  {\mathrm d\varepsilon}\right|_{0}
  \notag\\
  &\quad=
  \pmi_P(x,y)
  -\E_P\{\pmi_P(X,Y)\mid X=x\}
  -\E_P\{\pmi_P(X,Y)\mid Y=y\}
  +I(P).
  \label{eq:m2-score-term}
\end{align}

Combining Equations~\eqref{eq:mi-derivative},
\eqref{eq:g-decomposition}, \eqref{eq:m2-weight-term}, and
\eqref{eq:m2-score-term} yields
\begin{equation}
\boxed{
\begin{aligned}
g_P(x,y)
={}&\{\pmi_P(x,y)-I(P)\}^2-V(P)\\
&+2\bigl[
\pmi_P(x,y)
-\E_P\{\pmi_P(X,Y)\mid X=x\}\\
&\hspace{18mm}
-\E_P\{\pmi_P(X,Y)\mid Y=y\}
+I(P)
\bigr].
\end{aligned}
}
\label{eq:g-final}
\end{equation}
The first line is the change that would arise from the centred squared
pointwise-MI value. The remaining terms account for the row and column margins
that change with the selected cell.

\section{Sampling variance of the MI variance estimator}

Define the population variability of the observation-level sensitivities by
\begin{equation}
  \tau^2(P)=\Var_P\{g_P(X,Y)\}.
  \label{eq:tau-definition}
\end{equation}
The mean of an influence function is zero. Centring explicitly gives the same
variance if numerical substitution produces a small nonzero mean.

A first-order expansion of the variance estimator is
\begin{equation}
  \Vhat(P)-V(P)
  \approx
  \frac{1}{n_P}
  \sum_{a=1}^{n_P}g_P(Z_a^{(P)}).
  \label{eq:vhat-linearization}
\end{equation}
The observations are independent, so
\begin{align}
  \Var\{\Vhat(P)\}
  &\approx
  \frac{1}{n_P^2}
  \sum_{a=1}^{n_P}
  \Var_P\{g_P(Z_a^{(P)})\}\\
  &=\frac{\tau^2(P)}{n_P}.
  \label{eq:vhat-sampling-variance}
\end{align}
Substitution into Equation~\eqref{eq:component-df-moments} gives the
population component degree of freedom
\begin{equation}
  \boxed{
  \nu_P
  \approx
  \frac{2n_PV(P)^2}{\tau^2(P)}
  }.
  \label{eq:population-component-df}
\end{equation}
When \(\Vhat(P)\) is stable between samples, \(\tau^2(P)\) is small and the
component degrees of freedom are large. An unstable variance estimate produces
smaller degrees of freedom and a heavier-tailed final reference.

\section{Table-based component degrees of freedom}

The observed table replaces every population quantity in Equation
\eqref{eq:g-final} by its empirical counterpart. For each occupied cell,
\begin{equation}
\begin{aligned}
\widehat g_P(i,j)
={}&\{\widehat\pmi_P(i,j)-\Ihat(P)\}^2-\Vhat(P)\\
&+2\left[
\widehat\pmi_P(i,j)
-\frac{\sum_{j'}\widehat p_{ij'}\widehat\pmi_P(i,j')}
       {\widehat p_{i+}}
-\frac{\sum_{i'}\widehat p_{i'j}\widehat\pmi_P(i',j)}
       {\widehat p_{+j}}
+\Ihat(P)
\right].
\end{aligned}
\label{eq:g-hat}
\end{equation}
The empirical mean and variance of these cell sensitivities are
\begin{align}
  \overline g_P
  &=\sum_{i,j}\widehat p_{ij}\widehat g_P(i,j),\\
  \widehat\tau^2(P)
  &=\sum_{i,j}\widehat p_{ij}
  \{\widehat g_P(i,j)-\overline g_P\}^2.
  \label{eq:tau-hat}
\end{align}
The estimated component degree of freedom is
\begin{equation}
  \boxed{
  \widehat\nu_P
  =\frac{2n_P\Vhat(P)^2}{\widehat\tau^2(P)}
  }.
  \label{eq:nu-p-hat}
\end{equation}
Repeating the same calculation for the second table gives
\(\widehat\nu_Q\).

\section{Combined degrees of freedom and final test}

Substituting the two variance contributions and component degrees of freedom
into Equation~\eqref{eq:ws-general} gives
\begin{equation}
  \boxed{
  \widehat\nu_{\mathrm{expanded}}
  =
  \frac{
  \left\{\Vhat(P)/n_P+\Vhat(Q)/n_Q\right\}^2
  }{
  \left\{\Vhat(P)/n_P\right\}^2/\widehat\nu_P
  +
  \left\{\Vhat(Q)/n_Q\right\}^2/\widehat\nu_Q
  }
  }.
  \label{eq:expanded-df}
\end{equation}
The expanded two-sided p-value is
\begin{equation}
  p_{\mathrm{expanded}}
  =2\left\{1-F_{t_{\widehat\nu_{\mathrm{expanded}}}}(|T|)\right\},
  \label{eq:expanded-p}
\end{equation}
and the corresponding confidence interval is
\begin{equation}
  \Deltahat
  \ \pm\
  t_{\widehat\nu_{\mathrm{expanded}},1-\alpha/2}\SEhat.
  \label{eq:expanded-ci}
\end{equation}

The expanded calculation does not change \(\Deltahat\), \(\SEhat\), or
\(T\). It changes only the quantile used to evaluate that statistic.

\section{Worked count-table example}

Consider the two observed \(3\times3\) tables
\begin{equation}
  N^P=
  \begin{pmatrix}
  35&10&5\\
  10&20&5\\
  5&5&10
  \end{pmatrix},
  \qquad
  N^Q=
  \begin{pmatrix}
  30&5&5\\
  5&30&5\\
  5&5&10
  \end{pmatrix}.
  \label{eq:worked-tables}
\end{equation}
Their totals are \(n_P=105\) and \(n_Q=100\). The plug-in MI estimates are
0.13693 and 0.25848 nats. With \(d=(3-1)(3-1)=4\), the corrected difference is
\begin{equation}
  \Deltahat
  =\left(0.13693-\frac{4}{2(105)}\right)
  -\left(0.25848-\frac{4}{2(100)}\right)
  =-0.12059.
\end{equation}

The empirical influence variances are
\(\Vhat(P)=0.27189\) and \(\Vhat(Q)=0.43371\). Hence
\begin{equation}
  \SEhat
  =\left(\frac{0.27189}{105}+\frac{0.43371}{100}\right)^{1/2}
  =0.08323,
  \qquad
  T=-1.44898.
\end{equation}
Normal Wald gives \(p=0.14734\). Simple Welch uses total effective degrees of
freedom 188.53 and gives \(p=0.14900\).

For expanded Welch, the table-based sensitivity calculation gives component
degrees of freedom
\begin{equation}
  \widehat\nu_P=16.65,
  \qquad
  \widehat\nu_Q=45.86.
\end{equation}
Equation~\eqref{eq:expanded-df} combines them into
\(\widehat\nu_{\mathrm{expanded}}=59.03\), giving \(p=0.15264\). The three
methods reach the same substantive conclusion for these tables, but the
example shows where the expanded reference differs: it regards the two MI
variance estimates as less stable than the generic \(n-1\) assignment does.

\section{Operational algorithm and complexity}

Table~\ref{tab:expanded-algorithm} summarises one complete comparison.

\begin{table}[htbp]
\centering
\caption{Operational expanded Welch--Satterthwaite algorithm.}
\label{tab:expanded-algorithm}
\begin{tabularx}{\textwidth}{@{}cX@{}}
\toprule
Step & Calculation \\
\midrule
1 & Validate the two aligned independent count tables and their declared alphabet dimensions. \\
2 & Compute empirical probabilities, pointwise MI, plug-in MI, and the leading bias-corrected difference. \\
3 & Compute \(\Vhat(P)\), \(\Vhat(Q)\), the shared standard error, and \(T\). \\
4 & Compute the row and column pointwise-MI means needed in \(\widehat g_P\) and \(\widehat g_Q\). \\
5 & Compute \(\widehat\tau^2(P)\), \(\widehat\tau^2(Q)\), and the two component degrees of freedom. \\
6 & Combine the components using Equation~\eqref{eq:expanded-df}. \\
7 & Return the Student p-value, confidence interval, effective degrees of freedom, and validity diagnostics. \\
\bottomrule
\end{tabularx}
\end{table}

Each empirical probability, pointwise-MI value, row mean, column mean, and
variance contribution is calculated by a fixed number of table scans. The
time complexity is therefore \(O(rc)\), and the working memory is also
\(O(rc)\). The method is deterministic once the two count tables are given.

\section{Theoretical status}

The influence-function calculations in Equations~\eqref{eq:mi-derivative}
and \eqref{eq:g-final} are directional derivatives of the corresponding
population functionals. Under fixed positive support, they provide the usual
first-order asymptotic linear representations.

The scaled chi-squared model in Equation~\eqref{eq:scaled-chi-working} and the
final Student reference are working approximations. They match the first two
moments of the estimated variance rather than providing an exact finite-sample
law. Under regular large-sample conditions,
\(\widehat\nu_P\), \(\widehat\nu_Q\), and
\(\widehat\nu_{\mathrm{expanded}}\) increase, so the Student reference
approaches the same normal limit as Wald inference.

The method can fail when the empirical table loses too much support, when the
estimated MI variance is zero, or when the estimated sensitivity variance is
not positive and finite. These cases are recorded as invalid calculations
rather than silently replaced by a different test. The validation in Chapters
\ref{chap:design} and \ref{chap:results} evaluates how often these conditions
occur and whether the finite-sample correction improves calibration.
